% Part: counterfactuals
% Chapter: minimal-change-semantics
% Section: introduction

\documentclass[../../../include/open-logic-section]{subfiles}

\begin{document}

\olfileid[hi]{cnt}{min}{int}

\olsection{परिचय}

स्टालनेकर और लुईस ने ``यदि माचिस की तीली रगड़ी जाती, तो वह जलती'' जैसे
प्रतितथ्यात्मक सशर्तों के विवरण प्रस्तावित किए। उनके विवरण ऐसे वाक्यों की
सत्यता-शर्तों को ठीक से समझने के प्रस्ताव थे। दोनों प्रस्तावों का मूल विचार
यह है: कोई प्रतितथ्यात्मक सशर्त सत्य है या नहीं, इसका मूल्यांकन करने के लिए
हमें उन संभव जगतों पर विचार करना होगा जो पूर्ववर्ती को सत्य बनाने के लिए
वास्तविक जगत से न्यूनतम रूप से भिन्न हों। यदि इन संभव जगतों में परिणामांश
सत्य है, तो प्रतितथ्यात्मक सत्य है। मसलन, मान लें मेरे हाथ में माचिस की तीली
और माचिस की डिबिया है। वास्तविक जगत में मैं केवल उन्हें देखता और सोचता हूँ
कि तीली रगड़ने पर क्या होगा। वास्तविक जगत से वह न्यूनतम परिवर्तन, जिसमें
मैं तीली रगड़ता हूँ, ऐसा है जहाँ मैं कार्य करने का निर्णय लेकर तीली रगड़ता
हूँ। वह न्यूनतम इसलिए है कि और कुछ नहीं बदलता: मैं साथ ही हवा में नहीं
कूदता, तीली रगड़ने से मेरे बाल भी नहीं जल उठते, मेरी उँगलियों की सारी शक्ति
अचानक समाप्त नहीं होती, उसी क्षण सुपर-सोकर से घात लगाकर मुझ पर पानी नहीं
डाला जाता, इत्यादि। उस वैकल्पिक संभावना में तीली जलती है। अतः यह सत्य है
कि यदि मैं तीली रगड़ता, तो वह जलती।

इस सहज विवरण को प्रतितथ्यात्मक तर्कशास्त्र के औपचारिक अर्थविज्ञान के साथ
जोड़ा जा सकता है। लुईस ने प्रतितथ्यात्मक के लिए ``$\boxright$'' संकेत और
स्टालनेकर ने~``$>$'' संकेत प्रस्तुत किया। हम~$\cif$ का उपयोग करेंगे और उसे
प्रतिज्ञप्तिक तर्कशास्त्र में द्विस्थानी संयोजक के रूप में जोड़ेंगे। अतः
$!A \lif !B$ रूप वाले !!{formula}ों के अतिरिक्त हमारे पास~$!A \cif !B$ रूप
वाले !!{formula} भी हैं। मोडल तर्कशास्त्र के संबंधात्मक अर्थविज्ञान की तरह
औपचारिक अर्थविज्ञान ऐसे प्रतिरूपों पर आधारित है जिनमें जगतों पर !!{formula}ों
का मूल्यांकन होता है; और $\mSat{M}{!A \cif !B}[w]$ को परिभाषित करने वाली
संतुष्टि-शर्त कुछ (अन्य) जगतों~$w'$ पर $\mSat{M}{!A}[w']$ और
$\mSat{M}{!B}[w']$ के पदों में दी जाती है। कौन-से $w'$? सहज रूप से~$w$ के
वे निकटतम जगत जिनके लिए~$\mSat{M}{!A}[w']$। इसके लिए ``निकटता'' का संबंध भी
प्रतिरूप में सम्मिलित करना होगा।

लुईस ने प्रतितथ्यात्मक स्थितियों को आलेखीय रूप से निरूपित करने की उपयोगी
विधि प्रस्तुत की। प्रत्येक संभव जगत अन्य जगतों वाले अंतर्निविष्ट गोलकों के
एक समुच्चय के केंद्र में होता है---हम इन गोलकों को संकेंद्रित वृत्तों के रूप
में खींचते हैं। दो गोलकों के बीच के जगत केंद्र वाले जगत से समान दूरी पर हैं;
भीतरी गोलक में आने वाले जगत अधिक निकट और बाहरी गोलक में आने वाले अधिक दूर
हैं।
\begin{center}
\begin{tikzpicture}[scale=.7]\small
  \spheresystem{5}
  \draw (0,0) node {$w$};
  \propositionintersect{0}{4}{40}{3.5}
  {\tikzset{proposition/.append={smooth,tension=1.4}}
  \proposition[shift={(1.6,-1)}]{90}{1}{30}{4}}
  \path (1.5,3) node[below] {$\formula{B}$};
  \path (3,0) node {$\formula{A}$};
\end{tikzpicture}
\end{center}
निकटतम $!A$-जगत वे जगत~$w'$ हैं जहाँ~$!A$ संतुष्ट होता है और जो केंद्रीय
जगत~$w$ के चारों ओर सबसे छोटे गोलक (धूसर क्षेत्र) में स्थित हैं। सहज रूप से,
~$w$ पर $!A \cif !B$ तभी संतुष्ट होता है जब सभी निकटतम $!A$-जगतों पर $!B$
सत्य हो।

\end{document}
